\begin{lemma}[1.6, fixed local factors]\label{checked-smooth-factors}
\lean{NilBohr.auxiliary_step_dvd_fixed_power,NilBohr.localGeneralizedVonNeumann_on_lcm_moduli,NilBohr.centeredSelectorComparison_on_lcm_moduli_of_obstructions}\leanok
On the smooth moduli with $\mathcal L(w)\mid q$ reached by the
density iteration, the centered count is $O(w^{-A})$ for a fixed
$A>0$, uniformly over the families and inputs of Lemma~4.4.
\end{lemma}
\begin{proof}[Proof architecture]
\proves{checked-smooth-factors}\uses{external-inputs,checked-difference,p:heisenberg-phase}\leanok
An inverse auxiliary denominator $r\le w$ divides $q$.
Each resulting step $rq^b$ divides the common chosen power
$q^{B+1}$. Use the fixed-step correlations in the spatial
differencing argument and apply the centered phase bound
after substituting the polynomial shifts on the parameter cells.
\end{proof}

\begin{lemma}[Ambient counting comparison]\label{checked-ambient-discrete}
\lean{NilBohr.ambientSelectorCount_eq_explicit,NilBohr.ambientSelectorCount_centered,NilBohr.ambientSelectorCount_sharp,NilBohr.heisenbergAmbientSelectorCount_sharp_of_literature}\leanok
For the shifts $n,n^2$, set $M=\lfloor\sqrt N\rfloor$ and let
$\mathcal C(F)$ be the Mangoldt weighted count with selector test
$F(Z(n))$. For each fixed Lipschitz $F$,
$\mathcal C(F)-\mu_F\mathcal C(1)=o(1)$ uniformly in $A\subset[N]$.
For each fixed open arc $J$ the same comparison holds with
$F=\one_J$ and $\mu_F=|J|$.
The smooth comparison uses Proposition~6.2 and Ter\"av\"ainen's
fixed-family estimate as inputs. The sharp comparison also uses
the prime number theorem.
\end{lemma}
\begin{proof}[Proof architecture]
\proves{checked-ambient-discrete}\uses{p:prime-phase,external-inputs,checked-ergodic}\leanok
Use the prime phase bound and Ter\"av\"ainen's fixed-family
estimate with parameter length $M$ and spatial scale $M^2$.
The change to the ambient normalization costs a factor $M^2/N\le1$.
For interval cutoffs use upper and lower Lipschitz approximants and
the uniform bound on the count with test $1$.
\end{proof}

\begin{lemma}[Ambient Gowers obstructions]\label{checked-ambient-norm}
\lean{NilBohr.ambient_centered_gowers_prescribed_cutoff,NilBohr.tendsto_ambient_centered_localized_mode,NilBohr.ambient_centered_gowers_lower_bound,NilBohr.irrational_quadratic_gowers_two_tendsto_zero,NilBohr.quadratic_obstruction_common_window,NilBohr.no_uniform_weighted_gowers_two}\leanok
For the paper's $\alpha,\beta$ and Lipschitz $F$, prescribe
$h\ne0$ with $\widehat F(h)\ne0$ and a nonzero nonnegative
Lipschitz cutoff $\chi$ supported on an interval of length $<1/2$
inside $(0,1)$. Then, for all sufficiently large $M$,
\[
 0<|\widehat F(h)|\int_{\T}\chi
 \ll \norm{\widehat\chi}_{\ell^1}\norm{F(Z(\cdot))-\mu_F}_{U^3[M]}.
\]
The implicit constant is universal. In particular a nonconstant
$F$ gives nondecay at every fixed order $S\ge3$.
For any irrational $\theta$, the weight $\ee(-\theta n^2)$
has $U^2[M]\to0$, yet has count $1/2$ for suitable $1$-bounded
inputs supported on $[2M^2]$ for the shifts $n,n^2$.
The quadratic example is checked from irrationality alone.
\end{lemma}
\begin{proof}[Proof architecture]
\proves{checked-ambient-norm}\uses{checked-separating-gowers,checked-separating-quadratic}\leanok
At $X=M$ and $q=1$, scalar equidistribution gives the localized
correlation with limit $\widehat F(h)\int\chi$. The seam cutoff
equals one on the prescribed support. Apply the checked localized
multiplier and the order-three bound for the mean. For the quadratic example the
order-two power is a sum of squared linear exponential sums
over interval overlaps. The geometric bound holds away from a
small set of directions. The spatial phases cancel the quadratic
weight exactly on the common supporting count.
\end{proof}

\begin{lemma}[Pointwise limits with Bohr weights]\label{checked-bohr-weights}
\lean{NilBohr.ae_bohrPrimeMode_tendsto,NilBohr.ae_continuousBohrPrimeAverage_tendsto,NilBohr.ae_primeAverage_congruence_tendsto,NilBohr.ae_primeAverage_reduced_congruence_tendsto,NilBohr.ae_simultaneousBohrIntervalAverage_tendsto,NilBohr.simultaneousBohrIntervalDensity_tendsto,NilBohr.ae_selectedSimultaneousBohrIntervalAverage_tendsto}\leanok
For the polynomial ergodic averages, fixed linear twists,
continuous weights of finitely many linear torus coordinates,
fixed congruence classes, and simultaneous interval cutoffs have
almost-everywhere limits at the unrestricted prime normalization.
For reduced congruence classes, and for interval conditions with
positive limiting relative density, the selected-prime normalization
also converges. The proof uses the unrestricted prime pointwise theorem
on the product system, the prime number theorem and Vinogradov.
Reduced congruence normalization uses the progression form of the
prime number theorem.
\end{lemma}
\begin{proof}[Proof architecture]
\proves{checked-bohr-weights}\uses{external-inputs,checked-prime-distribution}
\leanok
Insert a circle rotation in an input at a linear shift, adding that
shift with constant input when needed. Then use the unrestricted
pointwise theorem and Fubini. Finite Fourier sums treat periodic
weights. Uniform trigonometric approximation treats continuous
weights, and boundary strips plus exact rational-grid interpolation
treat intervals, with either endpoint convention.
\end{proof}

\begin{lemma}[The introductory Bohr eigenfunction example]\label{checked-bohr-eigenfunction}
\lean{NilBohr.symmetricBohr_selected_exponential_tendsto,NilBohr.bohr_rotation_eigenfunction_example}\leanok
Let $Tt=t+\theta$ on $\T$ with $\theta$ irrational, $f(t)=\ee(t)$,
$P(n)=n$, and $0<\varepsilon<1/2$. The unrestricted prime average
of $f(T^{p-1}t)$ tends to zero. The average over primes with
$\tn{\theta(p-1)}<\varepsilon$, divided by their count, tends to
$\ee(t)\sin(2\pi\varepsilon)/(2\pi\varepsilon)\ne0$ for every $t$.
\end{lemma}
\begin{proof}
\proves{checked-bohr-eigenfunction}\uses{checked-prime-distribution}\leanok
Write the summand as $\ee(t)\ee(\theta(p-1))$. Vinogradov's
equidistribution gives zero for the unrestricted character average.
For the restricted average, approximate the indicator of the
open symmetric arc from above and below by continuous functions.
The prime density tends to $2\varepsilon$, while the restricted
character sum divided by the unrestricted prime count tends to
$\int_{-\varepsilon}^{\varepsilon}\ee(u)\,du=\sin(2\pi\varepsilon)/\pi$.
The ratio gives the stated limit. It is nonzero because
$0<2\pi\varepsilon<\pi$.
\end{proof}
\noindent\textbf{Lean status.} The open-arc restriction, prime
normalization and two differing limits are checked using the
Vinogradov interface in Lemma~\ref{checked-prime-distribution}.

\begin{lemma}[Progressions in selected primes]\label{checked-prime-progressions}
\lean{NilBohr.fixedWindow_primeProgression_of_sparse,NilBohr.greenTaoRelativeDensity_of_sparse,NilBohr.shrinkingWindow_primeProgression}\leanok
For fixed $J$, positive limiting relative prime density implies a
nontrivial progression of every length $k\ge2$. For $N\ge100$ and
$k\ge4$, the selected primes through
$N$ contain such a progression if their relative density exceeds
$C_k\exp(-(\log\log\log N)^{c_k})$, by Ter\"av\"ainen--Wang
Theorem~1.1. The constants depend only on $k$ and the latter
implication tests the actual density with the arc at scale $N$.
\end{lemma}
\begin{proof}[Proof architecture]
\proves{checked-prime-progressions}\uses{t:pointwise,external-inputs}\leanok
For the first claim, apply the finite density bound as $N\to\infty$,
testing length $\max(k,4)$ and taking a prefix of length $k$.
In the finite claim apply the density bound to the selected prime set.
At $k=4$ its stronger upper bound $(\log\log N)^{-a}$ is at most
$e^{1/(4a)}\exp(-(\log\log\log N)^{1/2})$.
\end{proof}

\begin{lemma}[Restricted maximal domination]\label{checked-maximal-domination}
\lean{NilBohr.selectedPrimeAverage_domination,NilBohr.selectedPrimeMaximal_domination,NilBohr.eLpNorm_selectedPrimeMaximal_le,NilBohr.memLp_selectedErgodicPrimeMaximal,NilBohr.selectedErgodicPrimeMaximal_congr_ae}\leanok
A fixed window with a positive relative-prime density limit has a
constant $C>0$ so that for every summand sequence and every cutoff,
the modulus of its selected-prime average is bounded by $C$
times the unrestricted prime average of the moduli. The same
constant transfers an $L^p$ estimate for the positive unrestricted
maximal operator to the selected maximal operator, including
$p=\infty$.
\end{lemma}
\begin{proof}[Proof architecture]
\proves{checked-maximal-domination}\uses{t:pointwise}\leanok
Compare positive sums and use the ratio of prime counts whenever
the selected set is nonempty. The limiting density bounds the ratio
eventually, and the finitely many earlier cutoffs give a finite
additional constant. An empty selected set has zero average.
Take the supremum in the extended nonnegative reals. Monotonicity
of $L^p$ norms gives the conditional maximal estimate. Measurability
and invariance of the ergodic maxima under a.e. replacement follow
from the countable set of iterates and cutoffs.
\end{proof}
